\section{Оценка сложности YETI}\label{sec:cheap-analysis}
В анализе используется функция Ляпунова
\[
 \Phi=\mathcal A\bigl(f(x)-f^\ast\bigr)
       +\tfrac12\norm{v-x^\ast}^2\le R^2.
\]
Принятие точки, предложенной эвристикой, не увеличивает первое слагаемое.
Проекция на $\mathcal H_y$ уменьшает второе. При ускоренном обновлении
$\Phi$ не возрастает благодаря условию точности локального решателя.
Отсюда $f(x)-f^\ast\le R^2/\mathcal A$.

Число итераций, на которых регуляризация увеличивается, ограничено
благодаря убыванию $\Phi$. На остальных итерациях регуляризация
уменьшается вдвое или растёт вес $\mathcal A$. Отсюда следует оценка
числа итераций до выполнения условия $\mathcal A\ge2R^2/\varepsilon$.
Соответствующие оценки доказаны в леммах~\ref{lem:outer-potential}
и~\ref{lem:outer-count} приложения~\ref{app:cheap-proofs}.
Каждая итерация вместе с предшествующим ей эвристическим блоком
требует в среднем $O(n\log N)$ запросов при фиксированной истории
предыдущих итераций.

\begin{theorem}[верхняя оценка]\label{thm:cheap-upper}
При предположениях~\ref{asm:components}--\ref{asm:accuracy}
алгоритм~\ref{alg:cheap} почти наверное останавливается и возвращает
точку $\hat x$, для которой
\begin{equation}\label{eq:cheap-upper-final}
 \E[f(\hat x)-f^\ast]\le\varepsilon,
\end{equation}
а математическое ожидание общего числа $T$ обращений к оракулу
удовлетворяет
\begin{equation}\label{eq:cheap-upper-cost}
 \E T\lesssim T_0\log^2N.
\end{equation}
В $T$ включены все обращения к оракулу в эвристических блоках,
основных итерациях и проверках.
\end{theorem}

Доказательство приведено
в приложении~\ref{app:upper-total}. Вариант, в котором допустимое число запросов
эвристики зависит от уже выполненной работы, рассмотрен
в приложении~\ref{sec:budget}.
